% T24 sources: simple_litreview.tex lines 43--54; closest_prior_work_table.tex
\subsection{Synthesis and Bounded Empirical Gap}

Table~\ref{tab:closest-prior-work} shows that the article's distinction does
not reduce to adding a DAG or a shared resource. The closest partial analogues
are distributed across several research areas: RCPSP provides precedence
relations and resource-capacity constraints; common-server and reentrant
models provide repeated service and different machine-retention rules; HRI
provides a queue for, and switching by, a single operator; general-purpose MAS
provide evidence on the effect of coordination topology; MAGIS provides an
automated review-and-revision cycle; and CAID provides concurrently executed
branches, merging, and testing. Each of these elements is known separately.

Within the core corpus assembled under the search protocol, with a cutoff date
of 2026-08-09, no study was identified that jointly calibrated, on software
tasks, a task- and mode-specific coefficient $k$, autonomous and human-attention
intervals, a dependency DAG, actual concurrency, repeated visits to a single
human, local retention of an agent slot, integration and switching overhead,
and the elapsed-time makespan of an accepted result. This statement applies
only to the reviewed corpus and the search channels specified in the protocol;
it is not a universal claim about the entire literature.

CAID, the closest software system, observes concurrency, elapsed completion
time, and integration, but does not include a human or a no-AI condition. MAGIS
includes automated review and revision but does not establish actual
concurrency or human service. Common-server and reentrant models formalize
repeated service but do not map it to the software process; HRI quantifies
waiting associated with human attention but not software work; and
general-purpose MAS evaluate primarily
success and quality rather than the makespan of a human--agent cell.

The formal contribution of the present work is the joint representation of
these quantities and the verification of model properties on synthetic
scenarios. The article itself does not perform empirical calibration of a
real-world workflow. Collecting observations on software tasks completed to an
accepted result, and prospectively validating whether a jointly estimated M4
predicts elapsed-time makespan better than the simpler levels, remain future
work.

\begin{landscape}
\begin{table}[!ht]
\centering
\caption{Closest formulation classes and the scope of the present model. The
values ``yes,'' ``partial,'' and ``no'' indicate whether a property is present
in a formulation class, not empirical equivalence or novelty.}
\label{tab:closest-prior-work}
\setlength{\tabcolsep}{1.1pt}
\renewcommand{\arraystretch}{1.18}
\renewcommand{\textsuperscript}[1]{\ensuremath{^{\mathrm{#1}}}}
\scriptsize
\begin{tabular}{>{\raggedright\arraybackslash}p{2.85cm}>{\centering\arraybackslash}p{1.45cm}>{\centering\arraybackslash}p{1.9cm}>{\centering\arraybackslash}p{1.65cm}>{\centering\arraybackslash}p{1.7cm}>{\centering\arraybackslash}p{1.6cm}>{\centering\arraybackslash}p{1.45cm}>{\centering\arraybackslash}p{1.55cm}>{\centering\arraybackslash}p{1.95cm}>{\raggedright\arraybackslash}p{2.6cm}>{\raggedright\arraybackslash}p{4.15cm}}
\toprule
Formulation class & DAG & \shortstack{Agent\\streams} & \shortstack{Unit-\\capacity\\service} & \shortstack{Repeated\\service} & \shortstack{Slot\\retention} & \shortstack{Task-\\and mode-\\specific $k$} & \shortstack{$a_i/h_i$\\phase split} & \shortstack{Coordination\\and\\integration\\overhead} & \shortstack{Primary\\outcome} & \shortstack{Real-world\\software-process\\calibration} \\
\midrule
RCPSP and multi-mode RCPSP \cite{brucker1999resource}\textsuperscript{a}
& yes & partial & yes & partial & no & no & no & partial & makespan, cost, feasibility & no \\

Parallel machines with a common server \cite{hall2000parallel,kravchenko1997parallel,elidrissi2023loading,chakhlevitch2009reentrant}\textsuperscript{b}
& no & yes & yes & partial & partial & no & partial & no & makespan, algorithmic bounds & no \\

HRI: robots per operator and attention allocation \cite{olsen2004fanout,crandall2005validating,crandall2011computing,cummings2008predicting}\textsuperscript{c}
& no & yes & yes & yes & partial & no & partial & partial & throughput, mission performance & no \\

General-purpose MAS and centralized coordination \cite{kim2025scaling}\textsuperscript{d}
& no & partial & no & partial & no & no & no & partial & success, quality; sometimes cost & no \\

Coding agents: MAGIS and CAID \cite{tao2024magis,geng2026caid}\textsuperscript{e}
& partial & partial & no & partial & no & no & no & partial & success, quality; CAID also time and cost & no \\

Present work\textsuperscript{f}
& yes & yes & yes & yes & yes & yes & yes & yes & makespan of an accepted result & no \\
\bottomrule
\end{tabular}

\vspace{0.5em}
\begin{minipage}{0.97\linewidth}
\scriptsize
\textsuperscript{a} An RCPSP renewable resource need not represent identical
agent streams; repeated service and overhead require explicit operations, and
the standard formulation does not introduce automatic slot retention while
another resource is awaited.

\textsuperscript{b} Repeated service is present in loading/unloading and
reentrant variants, but not in setup-only formulations. Immediate unloading
may retain a machine, whereas the reentrant model with a remote server releases
it; machine-processing and service durations are not calibrated as agent and
human components of software work.

\textsuperscript{c} An autonomous system waits for the operator or loses
performance, but this is not equivalent to retaining a software-development
slot. Interaction and autonomous intervals are not software-specific
$a_i/h_i$ components; queues, loss of situational awareness, and reorientation
do not include code integration.

\textsuperscript{d} Sequential interdependence is measured, but no explicit
task DAG is constructed; the number of agents or roles does not establish that
their intervals overlap. A centralized coordinator is not equivalent to
observed unit-capacity service, repeated messages are not a human service cycle,
and differences in outcomes do not constitute calibration of M4 overhead.

\textsuperscript{e} CAID constructs a dependency-aware plan and exhibits
overlapping branches; these properties are not established for MAGIS or for
the class as a whole. Automated planning, quality assurance, and revision are
not human service; waiting for dependencies or merging is not equivalent to
local blocking. CAID observes merging, testing, time, and cost, but does not
calibrate a human--agent process. Evaluation on software task suites is not
empirical calibration of a real-world workflow with observed human-attention
phases.

\textsuperscript{f} The article uses synthetic scenarios; empirical calibration
and prospective validation remain future work.

\medskip
The article does not claim novelty for DAGs, makespan, a shared resource,
repeated service, return to the same worker, resource retention,
resource-ordering arcs, or the number of autonomous workers per operator.
\end{minipage}
\end{table}
\end{landscape}
\FloatBarrier
